\section{Существующие поисковые системы}

По корпусу уже можно искать встроенным поиском Википедии, собственным поиском film.ru
(форма на каждой странице сайта) и поиском Google или Яндекса с ограничением
\texttt{site:}. Значит, будет с чем сравнивать качество собственной системы.

Ниже --- реальные запросы к встроенному поиску Википедии.

\textbf{[ я требую продолжения банкета ]} --- точная цитата из фильма \enquote{Иван
Васильевич меняет профессию}. Выдача: \emph{Иван Васильевич (пьеса)}, \emph{Революция
1848 года во Франции}, \emph{Скорая помощь}, \emph{Тёмный рыцарь}, \emph{Магическая
битва}. Статьи о фильме нет вообще, со второй позиции --- шум: революция 1848 года
попала сюда по одному слову \enquote{банкет}. Слова запроса ищутся независимо, цитата
не распознаётся как цитата.

\textbf{[ фильм где мужик теряет память ]} --- описательный запрос. Выдача: \emph{Иди и
смотри}, \emph{Высота}, \emph{Идиот (роман)} и две биографии актёров. Ни одного
правильного ответа: поиск сцепился за слова \enquote{фильм} и \enquote{память}, смысл
описания прошёл мимо, разговорное \enquote{мужик} ни с чем не сопоставлено.

\textbf{[ кто играл шерлока холмса в советском кино ]} --- выдача: \emph{Приключения
Шерлока Холмса и доктора Ватсона}, \emph{Майкрофт Холмс}, \emph{Шерлок Холмс (сериал,
2013)}, \emph{Клюев Б.\,В.}, \emph{Доктор Ватсон}. Первый документ верный, но статьи о
Василии Ливанове в пятёрке нет, зато есть исполнитель роли Майкрофта. Слова \enquote{кто}
и \enquote{в} не отброшены как стоп-слова, а британский сериал 2013 года противоречит
уточнению \enquote{в советском}.

\textbf{[ фильм про зону и сталкера ]} --- выдача: \emph{Сталкер (фильм)},
\emph{Сталкер}, \emph{Пикник на обочине}, \emph{Зона отчуждения ЧАЭС},
\emph{S.T.A.L.K.E.R.} Здесь всё хорошо. Запрос такой же описательный, как второй, но в
нём есть редкое слово \enquote{сталкер}. Видно закономерность: на редких информативных
терминах поиск справляется, на частотных словах общего языка --- разваливается.

\pagebreak
